\section{Limits of the approach}\label{sec:limits}

The run shows that the pipeline works, and it also shows where it is weak. We list the limits in the order of the stages of Figure~\ref{fig:pipeline}.

\paragraph{The arches are classical and shallow.} All five arches are characterisations on \emph{presheaf} topoi, where an invariant reduces to a finite condition on a small category. Such conditions are easy to state and easy to check, which is why every output has a proof of a few lines. The depth of a generated theorem is bounded by the depth of its arch and of its deck \cite[slide 8]{caramello-mines}; to get deeper outputs one needs arches on genuine sites (for instance Caramello's site characterisations of geometric invariants \cite{caramello-tst}) and decks between non-trivially equivalent theories, which our registry does not contain.

\paragraph{The decks are chosen by hand.} The new outputs came from decks that had not been paired with these arches before: quivers, synchronizing automata, transformation monoids. Choosing them was the creative step, and nothing in the pipeline automates it. A registry of \q{structure $X$ is a presheaf on shape $C$} is easy to populate for presheaf topoi; for classifying topoi of theories it requires proving Morita equivalences, which is research.

\paragraph{Specialising can be hard.} Computing the condition on the shape was trivial here (Ore on $\B M$, topologies on $\B\mathbb N$), but the count of Grothendieck topologies on $\B(\mathbb N^2)$, or the Ore condition on a category presented by generators and relations, is already a real problem. The generator moves the difficulty, it does not remove it.

\paragraph{Unfolding and rewriting need a human.} The raw output is \q{$\mathrm{Paths}(R^{\op})$ satisfies \texttt{RightOreCondition}}; the statement a mathematician wants is \q{every vertex has at most one outgoing arrow}. This step is what the talk assigns to rewriting systems and language models \cite[slide 17]{caramello-mines}, and it took most of our time.

\paragraph{Lean checks the corollary, not the reading.} A Lean proof that $\mathrm{Paths}(R^{\op})$ satisfies Ore iff a condition holds is a guarantee about that statement, under the library's conventions (side of the action, direction of arrows, \texttt{SingleObj} composing as $f\gg g = g\cdot f$). Whether the informal sentence in the article says the same thing is not checked by Lean; we met two sign conventions that flip the side of a monoid action, and the $T_n$ example shows that the side matters.

\paragraph{The novelty filter is weak.} \q{Not found} is not \q{new}. Our search (\S\ref{sec:lean-lit}) covered the obvious sources, but elementary statements are often folklore, stated in passing or in exercises, and not indexed. Most outputs of the first run were in print (\S\ref{sec:calibration}); the proportion is a fair measure of how far the simplest decks are from new mathematics.

\paragraph{Most of the output space is uninteresting.} Pairing every deck with every arch produces many statements that are true, new in the narrow sense, and of no interest (the weak excluded middle for presheaves on a particular five-object category). The talk's warning about \q{weird} results applies. A useful generator needs a ranking of outputs, which is a matter of mathematical taste and is not formalised.

\paragraph{What the approach is good for.} Within these limits, the pipeline does something useful: it asks questions nobody had asked (does the weak excluded middle hold for automata? for the actions of a synchronizing automaton's monoid?), answers them with a guarantee (the arch is formally proved), and points to a direct proof. As a tool for teaching, or for exploring the boundary between two fields, this is already worth having.
